\chapter{Versuchsurchführung}\label{ch:experiments}
\section{Aufnahme der Testdaten}
Über vier Monate hinweg werden Trainingsdaten von unterschiedlichen Straßentypen gesammelt. Der Fokus liegt auf den häufigsten Arten von Belägen, die in üblichen europäischen Städten anzutreffen sind: 

\begin{itemize}
    \item \textbf{Asphalt}
    \item \textbf{Verbund} (Betonstein- bzw.\ Verbundpflaster)
    \item \textbf{Pflastersteine} (Naturstein- bzw.\ Kopfsteinpflaster)
    \item \textbf{Feldweg} (Unbefestigter Weg, Erde, Schotter)
\end{itemize}

Diese werden zusätzlich in Qualitätsstufen unterteilt, um zu testen, welche maximale Granularität bei der Klassifizierung möglich und brauchbar ist.

\begin{itemize}
    \item \textbf{Asphalt}
    \begin{itemize}
        \item \textbf{Gut:} Glatte Oberfläche, frisch asphaltiert, kaum messbare Erschütterungen.
        \item \textbf{Mittel:} Raue Oberfläche, leichte Risse, ausgebesserte oder geflickte Stellen.
        \item \textbf{Schlecht:} Größere Schäden, Risse, Schlaglöcher, Baumwurzeln, starke und abrupte vertikale Beschleunigungen.
    \end{itemize}
    
    \item \textbf{Verbund}
    \begin{itemize}
        \item \textbf{Gut:} Eben verlegte Steine, enge Fugen, geschlossene, flache Oberfläche.
        \item \textbf{Mittel:} Leichte Setzungen im Untergrund, Fugen sind deutlich beim Fahren spürbar.
        \item \textbf{Schlecht:} Starke Fugen, hochstehende Kanten oder fehlende Steine.
    \end{itemize}
    
    \item \textbf{Pflastersteine}
    \begin{itemize}
        \item \textbf{Gut:} Gerades historisches Pflaster, relativ flach, erzeugt ein konstantes Vibrationsmuster.
        \item \textbf{Mittel:} Kopfsteinpflaster mit intakten Fugen. Erzeugt durchgehend starke Vibrationen.
        \item \textbf{Schlecht:} Sehr grobes, unregelmäßiges Pflaster mit tiefen Fugen, lockeren Steinen oder tiefen Kuhlen. Erzeugt extreme Erschütterungen.
    \end{itemize}
    
    \item \textbf{Feldweg}
    \begin{itemize}
        \item \textbf{Gut:} Festgefahrener, trockener Erdboden oder sehr feiner, verdichteter Schotter.
        \item \textbf{Mittel:} Loser Schotter, leichte Unebenheiten, Waschbrettmuster, flache Pfützen oder Spurrillen.
        \item \textbf{Schlecht:} Grobe Steine, tiefe Matschlöcher, Wege in oder am Rand von Wäldern.
    \end{itemize}
    \item \textbf{Stillstand} (Wenn Geschwindigkeit zu niedrig für erkennbare Datenmuster ist.)
\end{itemize}